\apendice{Formulário Engajamento do Usuário}
\label{ap:forms-ues}
\newcommand{\escalalikert}[3]{%
\begin{table}[H]
\centering
\caption{#1}
\begin{tabular}{L C C C C C L}
\toprule
\textbf{#2} & 1 & 2 & 3 & 4 & 5 & \textbf{#3} \\
\bottomrule
\end{tabular}
\legend{Fonte: Dados da pesquisa (2026).}
\end{table}
}
\section*{Apresentação}
Este questionário tem como objetivo avaliar o engajamento do(a) usuário(a) durante a interação com a plataforma, com base na \textit{User Engagement Scale} (UES). O instrumento está organizado em quatro seções, cada uma associada a uma dimensão do engajamento: \textit{Atenção e Imersão}, \textit{Usabilidade Percebida}, \textit{Apelo Visual e Estético} e \textit{Recompensa e Valor da Experiência}.

\section*{Instruções}
Reflita sobre sua experiência ao usar o sistema. Para cada afirmação, indique o quanto ela descreve sua experiência durante o jogo que você acabou de realizar, utilizando a escala abaixo:

\begin{quote}
\noindent 1 = Discordo totalmente \quad 2 = Discordo \quad 3 = Neutro \quad 4 = Concordo \quad 5 = Concordo totalmente
\end{quote}

\section{Seção 1 -- Atenção e Imersão}
Pense no quanto você esteve mentalmente presente durante o jogo.

\subsection{Eu me perdi nesta experiência (FA-S.1)}
\textit{(Resposta obrigatória)}

\escalalikert{Item FA-S.1}{Discordo totalmente}{Concordo totalmente}

\subsection{O tempo que passei usando a plataforma simplesmente escorregou (FA-S.2)}
\textit{(Resposta obrigatória)}

\escalalikert{Item FA-S.2}{Discordo totalmente}{Concordo totalmente}

\subsection{Eu estava absorto nesta experiência (FA-S.3)}
\textit{(Resposta obrigatória)}

\escalalikert{Item FA-S.3}{Discordo totalmente}{Concordo totalmente}

\section{Seção 2 -- Usabilidade Percebida}
Pense em como foi interagir com a interface da plataforma. \textbf{Atenção:} nesta seção, pontuações mais altas indicam experiência mais negativa\footnote{Os itens desta seção possuem pontuação reversa (\textit{reverse-coded}): valores mais altos correspondem a uma percepção mais negativa de usabilidade.}.

\subsection{Eu me senti frustrado ao usar a plataforma (PU-S.1 -- item reverso)}
\textit{(Resposta obrigatória)}

\escalalikert{Item PU-S.1 (reverso)}{Discordo totalmente}{Concordo totalmente}

\subsection{Achei a plataforma confusa de usar (PU-S.2 -- item reverso)}
\textit{(Resposta obrigatória)}

\escalalikert{Item PU-S.2 (reverso)}{Discordo totalmente}{Concordo totalmente}

\subsection{Usar a plataforma foi desgastante (PU-S.3 -- item reverso)}
\textit{(Resposta obrigatória)}

\escalalikert{Item PU-S.3 (reverso)}{Discordo totalmente}{Concordo totalmente}

\section{Seção 3 -- Apelo Visual e Estético}
Pense no visual e na aparência da plataforma.

\subsection{A plataforma foi visualmente atraente (AE-S.1)}
\textit{(Resposta obrigatória)}

\escalalikert{Item AE-S.1}{Discordo totalmente}{Concordo totalmente}

\subsection{A plataforma teve um design esteticamente agradável (AE-S.2)}
\textit{(Resposta obrigatória)}

\escalalikert{Item AE-S.2}{Discordo totalmente}{Concordo totalmente}

\subsection{A plataforma agradou meus sentidos visuais (AE-S.3)}
\textit{(Resposta obrigatória)}

\escalalikert{Item AE-S.3}{Discordo totalmente}{Concordo totalmente}

\section{Seção 4 -- Recompensa e Valor da Experiência}
Pense no valor e satisfação que o jogo te proporcionou.

\subsection{Usar a plataforma valeu a pena (RW-S.1)}
\textit{(Resposta obrigatória)}

\escalalikert{Item RW-S.1}{Discordo totalmente}{Concordo totalmente}

\subsection{Minha experiência foi recompensadora (RW-S.2)}
\textit{(Resposta obrigatória)}

\escalalikert{Item RW-S.2}{Discordo totalmente}{Concordo totalmente}

\subsection{Senti interesse genuíno nesta experiência (RW-S.3)}
\textit{(Resposta obrigatória)}

\escalalikert{Item RW-S.3}{Discordo totalmente}{Concordo totalmente}
